
% ════════════════════════════════════════════
% 模型建立（5.3.3.模型建立.tex，三级子块）
% ════════════════════════════════════════════
\subsubsection{含违约惩罚的滚动时域调整模型}

本节建立滚动时域下的计划购电与调整购电模型，并对计划与调整间的偏差费用作分段线性刻画。

\textbf{步骤1：0:00 计划购电策略}

在 0:00，微网依据当日 0:00 发布的光伏预报制定全天计划购电策略，设计划购电功率为 $g^{\mathrm{p}}_k$，计划模型为：
\begin{equation}
    \min \; Z^{\mathrm{p}}=\sum_{k=0}^{143} g^{\mathrm{p}}_k\,p_k\,\Delta t
\end{equation}
\begin{equation}
    \mathrm{s.t.}\quad g^{\mathrm{p}}_k+S^{0}_k+d^{\mathrm{p}}_k-c^{\mathrm{p}}_k\geq L_k
\end{equation}
\begin{equation}
    E^{\mathrm{p}}_{k+1}=E^{\mathrm{p}}_k+\eta\,c^{\mathrm{p}}_k\,\Delta t-\frac{1}{\eta}\,d^{\mathrm{p}}_k\,\Delta t
\end{equation}
\begin{equation}
    E_{\min}\leq E^{\mathrm{p}}_k\leq E_{\max},\quad 0\leq c^{\mathrm{p}}_k,d^{\mathrm{p}}_k\leq P_{\max}
\end{equation}
其中，$S^{0}_k$ 为 0:00 预报功率，$E^{\mathrm{p}}_0$ 为当日初始储电量，$E^{\mathrm{p}}_{144}$ 自由。

\textbf{步骤2：调整购电的偏差费用}

在后续预报时刻 $\tau\in\{6,12,18\}$，微网可对尚未执行的时段调整购电。设调整后购电功率为 $g^{\mathrm{a}}_k$，计划购电量高于调整购电量的部分按交易电价的 $50\%$ 计违约费用，调整购电量高于计划购电量的部分按交易电价的 $1.5$ 倍计价，则第 $k$ 个间隔的购电费用为：
\begin{equation}
    C_k=\Big[\min\{g^{\mathrm{p}}_k,g^{\mathrm{a}}_k\}+0.5\max\{g^{\mathrm{p}}_k-g^{\mathrm{a}}_k,0\}+1.5\max\{g^{\mathrm{a}}_k-g^{\mathrm{p}}_k,0\}\Big]p_k\Delta t
    \label{eq:q3_Ck}
\end{equation}
其中第一项为\textbf{实际执行的购电量}（两者较小者），上式等价于分段形式：
\[
\frac{C_k}{p_k\Delta t}=
\begin{cases}
0.5\,g^{\mathrm{p}}_k+0.5\,g^{\mathrm{a}}_k, & g^{\mathrm{a}}_k\leq g^{\mathrm{p}}_k,\\[6pt]
1.5\,g^{\mathrm{a}}_k-0.5\,g^{\mathrm{p}}_k, & g^{\mathrm{a}}_k\geq g^{\mathrm{p}}_k.
\end{cases}
\]
下调购电时（$g^{\mathrm{a}}_k<g^{\mathrm{p}}_k$）左段即实际购电按 $1$ 倍、未执行差额按 $0.5$ 倍计违约费；上调购电时（$g^{\mathrm{a}}_k>g^{\mathrm{p}}_k$）右段即计划内按 $1$ 倍、超出部分按 $1.5$ 倍加价。故 $C_k$ 关于 $g^{\mathrm{a}}_k$ 分段线性凸，在 $g^{\mathrm{a}}_k=g^{\mathrm{p}}_k$ 处连续，两段斜率为 $0.5p_k\Delta t$ 与 $1.5p_k\Delta t$，图像如下。
\begin{center}
\begin{tikzpicture}[>=Stealth, x=1.8cm, y=0.95cm]
  \draw[->] (-0.15,0) -- (2.75,0) node[below] {\small $g^{\mathrm{a}}_k/g^{\mathrm{p}}_k$};
  \draw[->] (0,-0.35) -- (0,3.25) node[left] {\small $\dfrac{C_k}{p_k\Delta t\,g^{\mathrm{p}}_k}$};
  \draw[thick, blue] (0,0.5) -- (1,1) -- (2,2.5);
  \draw[dashed, gray] (1,0) -- (1,1);
  \fill[red] (1,1) circle (1.6pt);
  \node[above right, red] at (1,1) {$g^{\mathrm{a}}_k=g^{\mathrm{p}}_k$};
  \foreach \x in {0,1,2} \draw (\x,0.06) -- (\x,-0.06) node[below] {\small $\x$};
  \foreach \y/\l in {0.5/0.5,1/1,2/2,2.5/2.5} \draw (0.06,\y) -- (-0.06,\y) node[left] {\small $\l$};
  \node[blue] at (0.32,1.0) {斜率 $0.5$};
  \node[blue] at (1.68,2.12) {斜率 $1.5$};
\end{tikzpicture}

{\bfseries 分段购电费用随调整购电量的变化（$g^{\mathrm{p}}_k$ 固定）}
\end{center}

\textbf{费用口径说明}：若按计划购电量全额结算，则 $g^{\mathrm{a}}_k=0$ 时用户未购电量却仍须支付 $1.5\,g^{\mathrm{p}}_k p_k\Delta t$，与“违约费仅针对未执行部分”相悖。故本文按实际购电量结算：基础电费按 $\min\{g^{\mathrm{p}}_k,g^{\mathrm{a}}_k\}$ 计，未执行差额按 $0.5$ 倍、超出部分按 $1.5$ 倍计价，以保证费用闭合。

\textbf{步骤3：分段费用的上境线性化}

为保持线性规划结构，引入上境变量 $w_k$，令其不小于分段费用的两段线性上界：
\begin{equation}
    w_k\geq 0.5\,p_k\Delta t\,(g^{\mathrm{a}}_k+g^{\mathrm{p}}_k)
    \label{eq:q3_w1}
\end{equation}
\begin{equation}
    w_k\geq 1.5\,p_k\Delta t\,g^{\mathrm{a}}_k-0.5\,p_k\Delta t\,g^{\mathrm{p}}_k
    \label{eq:q3_w2}
\end{equation}
由于目标为最小化 $\sum_k w_k$，最优解中 $w_k$ 将精确等于两段上界的较大者，即 $w_k=C_k$。于是调整子问题在预报时刻 $\tau$ 覆盖的时段 $[\tau,24:00]$ 内最小化：
\begin{equation}
    \min \sum_{k=6\tau}^{143} w_k
\end{equation}
\begin{equation}
    \mathrm{s.t.}\quad g^{\mathrm{a}}_k+S^{\tau}_k+d^{\mathrm{a}}_k-c^{\mathrm{a}}_k\geq L_k
\end{equation}
\begin{equation}
    E^{\mathrm{a}}_{k+1}=E^{\mathrm{a}}_k+\eta\,c^{\mathrm{a}}_k\,\Delta t-\frac{1}{\eta}\,d^{\mathrm{a}}_k\,\Delta t
\end{equation}
连同容量、功率约束及式~\eqref{eq:q3_w1}~与式~\eqref{eq:q3_w2}~的上境约束，初始储电量 $E^{\mathrm{a}}_{6\tau}$ 取上一决策执行到 $\tau$ 时刻的储电量。

\textbf{步骤4：紧急购电费用与全年总费用}

设滚动调整后最终执行的购电功率为 $g^{\mathrm{a}}_k$，最终放电、充电功率分别为 $d^{\mathrm{a}}_k$、$c^{\mathrm{a}}_k$，实际光伏功率为 $S^{\mathrm{r}}_k$。紧急购电的本质是实际供电不足负载的缺口，即：
\begin{equation}
    \delta_k=\max\Big\{L_k-\big(g^{\mathrm{a}}_k+S^{\mathrm{r}}_k+d^{\mathrm{a}}_k-c^{\mathrm{a}}_k\big),\,0\Big\}
\end{equation}
该式表示：最终策略执行后微网供电仍不足负载的部分才需以 5 倍电价紧急购电。特别地，若最终预报功率 $S^{\mathrm{u}}_k$（由 0:00、6:00、12:00、18:00 各时段分段拼接）下功率平衡取等号，即 $g^{\mathrm{a}}_k+S^{\mathrm{u}}_k+d^{\mathrm{a}}_k-c^{\mathrm{a}}_k=L_k$，则上式退化为 $\delta_k=\max\{S^{\mathrm{u}}_k-S^{\mathrm{r}}_k,0\}$；一般情况下调整购电与充放电会部分弥补光伏预报偏差，故应按实际供电缺口而非光伏预报差核算。紧急购电费用为：
\begin{equation}
    Z^{\mathrm{em}}_k=5\,p_k\,\delta_k\,\Delta t
\end{equation}
全年总购电费用为计划购电费用、调整相关费用与紧急购电费用之和：
\begin{equation}
    \min \; Z=\sum_{d=0}^{364}\sum_{k=0}^{143}\Big(C_{d,k}+5\,p_k\,\delta_{d,k}\Delta t\Big)
\end{equation}
该模型将滚动决策分解为四个顺序求解的线性规划，其思想与模型预测控制中滚动时域优化框架一致\cite{parisio2014}：计划、调整、紧急三层购电及其费用被组织为统一表达，前一层的求解结果作为后一层的已知输入顺次传递。下一节即按这四个线性规划依次求解，并逐日汇总全年费用。
